%%%%%%%%%%%%%%%%%%%%%%%%%%%%%%%%%%%%%%%%%%%%%%%%%%%%%%%%%%%%%%%%%%
\section{Construction of the manufactured FSI solution}
\label{app:mms}
%%%%%%%%%%%%%%%%%%%%%%%%%%%%%%%%%%%%%%%%%%%%%%%%%%%%%%%%%%%%%%%%%%
% Source: the sympy derivation mms_fsi.py and its check logs (logs in
% paper-data/mms_fsi/checks/). The closed form below follows that script.

\todo{The derivation script (sympy/symengine), the C++ code generator, the independent checks of the generated code and the case generator are to be added to the repository (\texttt{paper-data/mms\_fsi/}) once their packaging is agreed; this appendix should then point to them.}

This appendix gives the closed form of the manufactured solution of Section~\ref{sec:mms_fsi}, the design decisions behind it, and the checks of the construction and of the generated code.
The aim is a solution that (i) satisfies the kinematic and dynamic interface conditions exactly, so that the coupling is tested without an interface source; (ii) has a finite-strain solid and an interface that moves normal and tangential to itself, so that the fluid domain and its mesh deform; (iii) has fluid and solid fields of physically consistent size, in particular a fluid wall shear of the order of $\mu_F$ times a velocity gradient; and (iv) is closed form, so that it can be evaluated to round-off at any point.
The parameters are given in Table~\ref{tab:mms_params}; all quantities are in consistent units.

\begin{table}[htbp]
	\centering
	\small
	\caption{Manufactured FSI solution: parameters (exact rationals in the derivation). The flat variant uses the same values with the interface height of Equation~\eqref{eq:mms_flat} and $c_\kappa = 0$.}
	\label{tab:mms_params}
	\begin{tabular}{lll}
		\toprule
		Parameter & Symbol & Value \\
		\midrule
		Channel length, height; wall thickness & $L$, $H$; $T$ & $2$, $1$; $1/2$ \\
		Angular frequency (period 1) & $\omega$ & $2\pi$ \\
		Normal interface amplitude; travelling-wave fraction & $a_\eta$; $b_\eta$ & $1/10$; $1/2$ \\
		Tangential stretch amplitude & $c_\kappa$ & $1/10$ \\
		Through-thickness shape factor & $\gamma$ & $1/4$ \\
		Interface shear amplitude & $e$ & $1/50$ \\
		Outlet flux amplitude; vortex amplitude & $Q_0$; $A_\Omega$ & $1$; $8$ \\
		Free pressure amplitude & $\pi_0$ & $1/2$ \\
		Fluid density, dynamic viscosity & $\rho_F$, $\mu_F$ & $1$, $1/100$ \\
		Solid density, Young's modulus, Poisson's ratio & $\rho_S$, $E_S$, $\nu_S$ & $1$, $10$, $3/10$ \\
		\bottomrule
	\end{tabular}
\end{table}

\paragraph{Time factors.}
All fields are multiplied by
\begin{equation}
	r(t) = \tfrac12\left(1 - \cos\omega t\right)
	\qquad\text{or}\qquad
	r_2(t) = \sin^2\omega t,
\end{equation}
which vanish with their first derivatives at $t = 0$.
The displacement, the solid velocity, the fluid velocity and the pressure are therefore zero at $t = 0$, and the calculation starts from rest on an undeformed mesh.

\paragraph{Solid first: the interface motion.}
The kinematic condition, $\bb{v}(\bb{x}^\Gamma(X,t),t) = \partial\bb{u}/\partial t\,(X,H,t)$, holds at the deformed position of the material interface point $X$, whereas the fluid is written in Eulerian coordinates.
The construction therefore starts from the solid, and requires the interface parametrisation $X \mapsto x^\Gamma(X,t)$ to be invertible in closed form.
The tangential interface motion is a M\"obius map of $[0,L]$ onto itself,
\begin{equation}
	x^\Gamma(X,t) = \frac{L\,q}{q + \kappa(t)\,(1-q)},
	\qquad q = \frac{X}{L},
	\qquad \kappa(t) = 1 + c_\kappa\,r_2(t),
\end{equation}
which fixes both ends and moves the interior points with a non-uniform tangential velocity; its inverse is rational as well,
\begin{equation}
	\Xi(x,t) = \frac{L\,\kappa\,\hat{x}}{1 - \hat{x} + \kappa\,\hat{x}},
	\qquad \hat{x} = \frac{x}{L}.
\end{equation}
The tangential interface displacement is $\xi(X,t) = x^\Gamma(X,t) - X$.

\paragraph{Eulerian interface height.}
The normal motion is prescribed in Eulerian form, as the channel height
\begin{equation}
	h(x,t) = H + \hat\eta(x,t),
	\qquad
	\hat\eta(x,t) = a_\eta\,\sin^2\!\left(\frac{\pi x}{L}\right) r(t) \left[1 + b_\eta \sin\!\left(\frac{2\pi x}{L} - \omega t\right)\right],
	\label{eq:mms_height}
\end{equation}
a pulsing wave with a travelling part, and the normal solid interface displacement follows by composition, $\eta(X,t) = \hat\eta(x^\Gamma(X,t),t)$.
Specifying $\eta$ in Lagrangian form instead makes the swept area $\int h\,\mathrm{d}x$, which the fluid flux below requires, non-elementary; in Eulerian form it is elementary.
The factor $\sin^2(\pi x/L)$ gives zero height and zero slope at the clamped ends, which the shear design below needs.
In the flat variant the interface stays flat and moves only vertically,
\begin{equation}
	\hat\eta(t) = a_\eta\,r(t)\left[1 + b_\eta \sin\omega t\right], \qquad \kappa = 1,
	\label{eq:mms_flat}
\end{equation}
and the solid ends carry the exact displacement instead of a clamp.

\paragraph{Bulk solid displacement.}
With $\hat y = (Y-H)/T$ and $g(Y) = (1-\hat y)(1+\gamma\hat y)$, so that $g = 1$ on the interface and $g = 0$ on the top,
\begin{align}
	u_X(X,Y,t) &= \xi(X,t)\,g(Y) + (Y-H)(1-\hat y)^2\left[m_0(X,t) + m_1(X,t)\right], \\
	u_Y(X,Y,t) &= \eta(X,t)\,g(Y),
\end{align}
where $m_0$ and $m_1$ change only the through-thickness shear on the interface (below).
The displacement vanishes on the top, $Y = H+T$, and, in the full variant, on the ends $X = 0$ and $L$, where $\xi$, $\eta$, $\partial\eta/\partial X$, $m_0$ and $m_1$ vanish; the interface is the only boundary of the solid with non-homogeneous data.
The solid is the compressible neo-Hookean law of the code in plane strain,
\begin{equation}
	\bb{\sigma}_S = \frac{\mu_S}{J}\,\mathrm{dev}\,\bar{\bb{b}} + \frac{K_S}{2}\,\frac{J^2-1}{J}\,\mathbf{I},
	\qquad
	\bar{\bb{b}} = J^{-2/3}\,\bb{F}\bb{F}^{\mathrm{T}},
	\qquad
	\bb{P} = J\,\bb{\sigma}_S\,\bb{F}^{-\mathrm{T}},
\end{equation}
with $\bb{F} = \mathbf{I} + \nabla_0\bb{u}$ ($F_{zz} = 1$), $J = \det\bb{F}$, $\mu_S = E_S/(2(1+\nu_S))$ and the three-dimensional bulk modulus $K_S = \lambda_S + \tfrac23\mu_S$, $\lambda_S = \nu_S E_S/((1+\nu_S)(1-2\nu_S))$.

\paragraph{The neo-Hookean shear identity.}
For an arbitrary solid displacement the interface shear traction is of order $\mu_S$ times the strain, i.e.\ $O(1)$ here; matching it with a fluid of viscosity $\mu_F = 0.01$ needs fluid shear rates of order $10^2$, and gave fluid velocities of about 30 in a first attempt.
For the law above, let $\bb{a} = \partial\bb{x}/\partial X$ and $\bb{b} = \partial\bb{x}/\partial Y$ be the deformed material tangent and through-thickness fibre on the interface.
The volumetric part of $\bb{\sigma}_S$ and the isotropic part of $\mathrm{dev}\,\bar{\bb{b}}$ give no shear, and since the deformed normal is parallel to $\bb{F}^{-\mathrm{T}}\bb{N}$ and $|\bb{F}^{-\mathrm{T}}\bb{N}| = |\bb{a}|/J$ in two dimensions, the shear traction on the interface is
\begin{equation}
	\left|\bb{\tau}\cdot\bb{\sigma}_S\,\bb{n}\right| = \mu_S\,J^{-2/3}\,\frac{|\bb{a}\cdot\bb{b}|}{|\bb{a}|^2},
\end{equation}
with $\bb{\tau}$ and $\bb{n}$ the unit tangent and normal of the deformed interface.
It vanishes if and only if the deformed fibre stays orthogonal to the deformed interface.
Since $\bb{b} = (\xi g_0' + m_0 + m_1,\ 1 + \eta g_0')$ and $\bb{a} = (1 + \partial\xi/\partial X,\ \partial\eta/\partial X)$ on $Y = H$, with $g_0' = \mathrm{d}g/\mathrm{d}Y\,(H) = (\gamma-1)/T$, the choice
\begin{equation}
	m_0 = -\frac{\partial\eta/\partial X\,(1 + \eta\,g_0')}{1 + \partial\xi/\partial X} - \xi\,g_0'
\end{equation}
makes $\bb{a}\cdot\bb{b} = 0$, and
\begin{equation}
	m_1 = e\,\sin\!\left(\frac{2\pi X}{L}\right) r(t)
\end{equation}
adds a small, smooth shear of known size, with a peak of $0.07$ against a peak normal traction of $2.6$, so that the fluid wall shear stays of order $\mu_F$ times a velocity gradient.
The identity is specific to this law: for St Venant--Kirchhoff the interface shear is quadratic in $\bb{b}$, and another device would be needed.

\paragraph{Divergence-free velocity on the deforming channel.}
The fluid velocity is the curl of a stream function written in the stretched coordinate $s = y/h(x,t) \in [0,1]$,
\begin{equation}
	\bb{v} = \left(\frac{\partial\psi}{\partial y},\ -\frac{\partial\psi}{\partial x}\right),
	\qquad
	\psi(x,y,t) = \Psi\!\left(x, \frac{y}{h(x,t)}, t\right),
\end{equation}
so that $\nabla\cdot\bb{v} = 0$ identically and the pressure-correction algorithms need no mass source.
With the interface velocity at the Eulerian abscissa, $\bb{V}(x,t) = \partial\bb{u}/\partial t\,(\Xi(x,t), H, t)$,
\begin{equation}
	\Psi(x,s,t) = Q(x,t)\,B(s) + h(x,t)\,V_x(x,t)\,C(s) + W_2(s)\,\Phi(x,t) + W_3(s)\,\Omega(x,s,t),
	\label{eq:mms_psi}
\end{equation}
with the Hermite-type functions $B = 3s^2 - 2s^3$, $C = s^3 - s^2$, $W_2 = s^2(1-s)^2$ and $W_3 = s^2(1-s)^3$.
On the wall $s = 0$ every term vanishes with its $s$-derivative (no-slip).
On the interface $s = 1$, $\psi = Q$ and $\partial\psi/\partial y = V_x$, so that the tangential velocity matches, and $v_y = -\partial Q/\partial x + (\partial h/\partial x)\,V_x$; the normal velocity therefore matches, $v_y = \partial h/\partial t + (\partial h/\partial x)\,V_x = V_y$ (the rate of change of the height following the sliding interface point), provided that $\partial Q/\partial x = -\partial h/\partial t$.
$Q$ is the volume flux through the section at $x$, and is chosen as
\begin{equation}
	Q(x,t) = Q_0\,r(t) + \frac{\partial}{\partial t}\int_x^L \hat\eta(x',t)\,\mathrm{d}x',
\end{equation}
which prescribes a non-negative outlet flux $Q(L,t) = Q_0\,r(t)$ and lets the inlet absorb the change of volume of the channel; the velocity-gradient condition at the outlet then never meets backflow, which makes such an outlet unstable.
The last term of Equation~\eqref{eq:mms_psi} is a travelling vortex train,
\begin{equation}
	\Omega = A_\Omega\,r(t)\,\sin^2\!\left(\frac{\pi x}{L}\right)\sin\!\left(\frac{2\pi x}{L} - \omega t\right)\left(1 + \tfrac12\cos\pi s\right),
\end{equation}
faded out at both ends so that it does not reverse the outlet flow; its factor $W_3$ vanishes with its first two derivatives at $s = 1$, so it changes neither the interface velocity nor the interface traction.

\paragraph{Exact traction continuity without an interface source.}
The unit normal and tangent of the interface $y = h(x,t)$ are $\bb{n} = (-\partial h/\partial x, 1)/\sqrt{1+(\partial h/\partial x)^2}$ and $\bb{\tau} = (1, \partial h/\partial x)/\sqrt{1+(\partial h/\partial x)^2}$; $\bb{n}$ agrees with the Nanson normal $\bb{F}^{-\mathrm{T}}\bb{N}$ of the solid.
The fluid traction $\bb{\sigma}_F\bb{n} = -p\,\bb{n} + 2\mu_F\,\bb{D}(\bb{v})\,\bb{n}$ is linear in two functions that are still free on the interface: the stream-function term $\Phi$ and the interface pressure $p_\Gamma$.
Since $W_2$ has zero value and slope at $s = 1$, $\Phi$ and its derivatives do not change the interface velocity, and $\Phi$ enters the interface shear rate only through $W_2''(1) = 2$; the tangential balance is therefore algebraic in $\Phi$,
\begin{equation}
	\bb{\tau}\cdot 2\mu_F\bb{D}\,\bb{n} = c_0 + c_1\Phi = \bb{\tau}\cdot\bb{\sigma}_S\,\bb{n},
	\qquad
	c_1 = \frac{2\mu_F\left(1 + (\partial h/\partial x)^2\right)}{h^2},
\end{equation}
where $c_0$ is the shear of the other terms of Equation~\eqref{eq:mms_psi} and $\bb{\sigma}_S$ is evaluated at $(\Xi(x,t), H, t)$; hence $\Phi = (\bb{\tau}\cdot\bb{\sigma}_S\bb{n} - c_0)/c_1$, in closed form.
The normal balance gives the interface pressure, and the pressure field is completed with a free term that vanishes on the interface,
\begin{equation}
	p_\Gamma = 2\mu_F\,\bb{n}\cdot\bb{D}\,\bb{n} - \bb{n}\cdot\bb{\sigma}_S\,\bb{n},
	\qquad
	p(x,y,t) = p_\Gamma(x,t) + (1-s)\,\pi_0\,r(t)\cos\!\left(\frac{\pi x}{L}\right).
\end{equation}
Both interface conditions are thus satisfied exactly, and no interface source is needed.
An interface source would have required a new term, implemented consistently, in each interface treatment (the quasi-monolithic traction transfer, the Neumann condition of the Dirichlet--Neumann coupling and the Robin condition), and would have bypassed part of the traction transfer under test.

\paragraph{Sources.}
The momentum sources are the residuals of the exact fields,
\begin{equation}
	\bb{f} = \rho_F\left(\frac{\partial\bb{v}}{\partial t} + (\nabla\bb{v})\,\bb{v}\right) - \nabla\cdot\bb{\sigma}_F,
	\qquad
	\bb{b}_0 = \rho_S\,\frac{\partial^2\bb{u}}{\partial t^2} - \nabla_0\cdot\bb{P},
\end{equation}
per unit current volume of fluid and per unit reference volume of solid; both enter the equations as body forces.
The Eulerian fluid source applies unchanged on the moving mesh, because $\partial\bb{v}/\partial t|_{\chi} + (\nabla\bb{v})(\bb{v} - \bb{w}) = \partial\bb{v}/\partial t|_{\bb{x}} + (\nabla\bb{v})\,\bb{v}$ for any mesh velocity $\bb{w}$.
The fluid code uses the kinematic pressure, so the fluid source is applied per unit mass; it is evaluated at the current cell centres at the new time level, and the solid source at the reference cell centres.
The sources are of the size of the physical terms: in the fluid $|\bb{f}| \approx 10$, against $|\partial\bb{v}/\partial t| \approx 9$, $|(\nabla\bb{v})\bb{v}| \approx 6$ and $|\nabla p| \approx 8$; in the solid $|\bb{b}_0| \approx 10$, against $|\rho_S\,\partial^2\bb{u}/\partial t^2| \approx 5$ and $|\nabla_0\cdot\bb{P}| \approx 11$.

\paragraph{Boundary data.}
All boundary data are exact.
The fluid has the exact velocity at the inlet, the exact normal velocity gradient and the exact pressure at the outlet (which fixes the pressure level), and no-slip on the bottom wall.
The pressure at the inlet and on the bottom wall has a flux-consistent normal gradient in the partitioned fluid, and the exact or momentum-balance normal gradient in the quasi-monolithic fluid.
The fluid interface conditions are those of the coupling.
The solid is clamped on its top and, in the full variant, on its ends; in the flat variant the ends carry the exact displacement.

\paragraph{Mesh motion.}
The exact solution does not depend on the motion of the fluid mesh, but the discrete error does.
In the full variant the ends of the fluid domain do not move, because $\hat\eta$ vanishes there and the M\"obius map fixes the end points of the interface; the interior mesh follows the interface by a displacement Laplacian with quadratic inverse-distance diffusivity.
This moves the cells next to the interface almost vertically, so that they become non-orthogonal to the tilting interface by up to $13^\circ$; a stress-based mesh motion halves this to about $6.5^\circ$, with similar errors.
This non-orthogonality is what the flat variant removes.
In the flat variant the points of the inlet and outlet boundaries are given the exact linear stretching $y \mapsto y\,h(t)/H$; with a slip condition there, the corner point lags behind the interface and folds the corner cell.

\paragraph{Checks of the construction.}
The fields are built in sympy and differentiated in symengine; eleven properties are checked, symbolically or at 40 significant digits (133-bit MPFR) at random points: the closed-form inverse $\Xi(x^\Gamma(X,t),t) = X$ (symbolic); the identity $\nabla\cdot\nabla\times\psi = 0$ for a generic $\psi$ (symbolic); that the top of the fluid domain is the deformed interface, $h(x^\Gamma) = y^\Gamma$; $\nabla\cdot\bb{v} = 0$; the kinematic condition $\bb{v}(\bb{x}^\Gamma) = \partial\bb{u}/\partial t$; no-slip on $y = 0$; traction continuity $\bb{\sigma}_F\bb{n} = \bb{\sigma}_S\bb{n}$ with the Nanson normal; zero solid interface shear when $m_1 = 0$; the coefficient $c_1$; the clamps; and the zero initial state.
The largest residuals are $2\times10^{-38}$ (kinematic condition) and $1.4\times10^{-38}$ (traction), against an interface traction of order one.
The checks pass for the full and the flat variants, and for a set with $\rho_F = 2$ and $\rho_S = 3$ that checks the density scaling.

\paragraph{Checks of the generated code.}
The closed-form expressions are long, and are evaluated in double precision by generated C++ (common-subexpression elimination; about 3200 temporaries for the fluid fields).
A separate program checks the generated code independently of the symbolic derivation, at 200 random points: fourth-order finite differences of the generated velocity and pressure reproduce the generated fluid source (to $4.9\times10^{-10}$, against $|\bb{f}| \approx 9$) and the generated gradients, and give a zero divergence ($3.8\times10^{-11}$); the neo-Hookean law is implemented a second time, line by line from the source of the code, and applied to finite differences of the generated displacement, giving the stress ($3.6\times10^{-10}$) and the solid source $\rho_S\,\partial^2\bb{u}/\partial t^2 - \nabla_0\cdot\bb{P}$ ($2.2\times10^{-7}$, the truncation error of the nested differences, against $|\bb{b}_0| \approx 8$); the generated solid velocity equals the finite-difference rate of the displacement; the interface normal from the finite-difference tangent agrees, and $h(x^\Gamma) = y^\Gamma$; the fluid velocity at the interface equals the solid velocity and the rate of the interface position; and the fluid and solid tractions agree ($6.7\times10^{-15}$).
The generated values also agree with the 40-digit symbolic values to $2.3\times10^{-14}$ (relative).
The program exits with a non-zero status if any check fails; all checks pass for the three parameter sets above (\texttt{paper-data/mms\_fsi/checks/}).
